\documentclass[conference]{IEEEtran}
\IEEEoverridecommandlockouts
% The preceding line is only needed to identify funding in the first footnote. If that is unneeded, please comment it out.
%Template version as of 6/27/2024

\usepackage{cite}
\usepackage{amsmath,amssymb,amsfonts}
\usepackage{algorithmic}
\usepackage{graphicx}
\usepackage{textcomp}
\usepackage{xcolor}
\def\BibTeX{{\rm B\kern-.05em{\sc i\kern-.025em b}\kern-.08em
    T\kern-.1667em\lower.7ex\hbox{E}\kern-.125emX}}
\begin{document}

\title{L-RAPID: Lightweight Reversible And Privacy-preserving IIoT Data-perturbation system
\thanks{Identify applicable funding agency here. If none, delete this.}
}

\author{\IEEEauthorblockN{1\textsuperscript{st} Given Name Surname}
\IEEEauthorblockA{\textit{dept. name of organization (of Aff.)} \\
\textit{name of organization (of Aff.)}\\
City, Country \\
email address or ORCID}
\and
\IEEEauthorblockN{2\textsuperscript{nd} Given Name Surname}
\IEEEauthorblockA{\textit{dept. name of organization (of Aff.)} \\
\textit{name of organization (of Aff.)}\\
City, Country \\
email address or ORCID}
\and
\IEEEauthorblockN{3\textsuperscript{rd} Given Name Surname}
\IEEEauthorblockA{\textit{dept. name of organization (of Aff.)} \\
\textit{name of organization (of Aff.)}\\
City, Country \\
email address or ORCID}
\and
\IEEEauthorblockN{4\textsuperscript{th} Given Name Surname}
\IEEEauthorblockA{\textit{dept. name of organization (of Aff.)} \\
\textit{name of organization (of Aff.)}\\
City, Country \\
email address or ORCID}
\and
\IEEEauthorblockN{5\textsuperscript{th} Given Name Surname}
\IEEEauthorblockA{\textit{dept. name of organization (of Aff.)} \\
\textit{name of organization (of Aff.)}\\
City, Country \\
email address or ORCID}
\and
\IEEEauthorblockN{6\textsuperscript{th} Given Name Surname}
\IEEEauthorblockA{\textit{dept. name of organization (of Aff.)} \\
\textit{name of organization (of Aff.)}\\
City, Country \\
email address or ORCID}
}

\maketitle

\begin{abstract}
IIoT(Industrial Internet of Things) systems constantly stream telemetry data from sensors and machinery to the cloud, posing serious privacy and security challenges. Conventional encryption methods are computationally too expensive for legacy devices involving complex mathematical operations like fiestel structures, discrete logarithm, etc., while common noise-based perturbation methods irreversibly distort data. We introduce and evaluate a lightweight, fully reversible data perturbation pipeline for IIoT telemetry, contributing a composite reversibility privacy-key-space scoring framework for comparing perturbation strategies under realistic legacy-device constraints. To assess the balance between privacy and utility, a dataset comprising 31,106 IIoT intrusion detection samples covering five distinct attack categories was employed. After training a classifier that achieved a 98.47% binary classification accuracy on original data, we evaluated four different perturbation techniques: Gaussian noise, Laplacian noise, XOR 16-bit, and XOR 64-bit against both a static and an adaptive adversary.
\end{abstract}

\begin{IEEEkeywords}
component, formatting, style, styling, insert.
\end{IEEEkeywords}

\section{Introduction}
IIoT equipment produces valuable operational telemetry data, the capture of which risks revealing commercially sensitive information about manufacturing operations and system design. Complete cryptographic shielding is too expensive for legacy equipment, while the noise addition approach may introduce permanently distorted signals. Although XOR masking, Gaussian noise and Laplacian noise are each individually well-established techniques, prior work has not, to our knowledge, evaluated them side-by-side under a unified reversibility, privacy and key-space scoring framework on a realistic IIoT dataset nor tested them against an adversary that retrains directly on the perturbed signal. This work bridges the gap by evaluating four approaches on a realistic IIoT dataset with 31,106 samples (15,000 normal, 16,106 attack) and designing a rigorous composite scoring system for evaluation. 

\section{Methodology}

We used the Modbus telemetry subset of the publicly available ToN IoT dataset, comprising 31,106 samples, labeled as normal 15,000 (48 \%) and attack 16,106 (52 \%) across five attack categories: Injection, Backdoor, scanning, XSS, and Password. The four Modbus function-code registers (FC1-FC4) are stationary, with low variance ratios ranging from 0.56 to 0.58, indicating that the raw signal is information-rich and difficult to model externally.

The Random Forest Classifier achieved 98.47\% (weighted F1: 0.9847) on a binary detection of normal versus attack traffic, with 16 (of 3000) normal samples classified as attacks and 79 (of 3222) attack samples missed on the binary task. This provides evidence that the raw IIoT telemetry network has a highly exploitable discriminative structure in its features and simultaneously justifies privacy-preserving perturbation before transmission. 

The study evaluated four distinct perturbation strategies against two adversary models:a static adversary, represented by a classifier trained on the original data and evaluated directly on perturbed samples and a stronger adaptive adversary who retrains a fresh Random Forest classifier directly on the perturbed values using an 80/20 stratified train-test split modeling an attacker with access to ground-truth labels for a substantial portion of intercepted traffic XOR masking, uses a bitwise exclusive-OR (XOR) operation to combine data values with an undisclosed secret key (describe derivation e.g., derived per-session from a PRNG seeded). Because this deterministic process preserves the data without introducing randomness, it remains entirely reversible. The 16-bit approach uses a 65,536 key space, while the 64-bit approach stretches it to 264. The testing phase also included additive noise utilizing Gaussian and Laplacian distributions. To guarantee consistency across these randomized methods, a uniform random seed was employed. The Laplacian distribution's "fatter" tails caused perturbed data points to be displaced significantly farther from their original positions than in the Gaussian distribution at equivalent power levels, thereby increasing the difficulty of statistical unperturbation for an adversary. While all four algorithms successfully achieved 100% exact-match restoration, they simultaneously hindered attackers' efforts, typically lowering accuracy to 47-51%. 50% adversarial accuracy in binary classification implies the attacker has no knowledge from which to infer the results, and thus the probability of each outcome is 0.5.

XOR 64-bit is the most advantageous perturbation technique among the methods evaluated. Its computational efficiency is high (can add real timing/cycle comparison), requiring only a single CPU instruction per value, and it allows for deterministic data reconstruction without the need to store a seed. Furthermore, the massive keyspace makes brute-force recovery attempts impractical.Gaussian and Laplacian perturbation also achieved full reversibility and comparable privacy protection, but both rely on a single shared seed for both encryption and decryption, a smaller and less formally characterized secret space than XOR-64’s explicit key and should be paired with stronger seed management practices before production use. The uniformity of low attack accuracy seen by each method suggests that what is necessary to achieve secure communication is not a particular method of perturbation, but, in general, to eradicate the statistical structure of the transmitted signal, in whatever form, or by whichever of the introduced methods. Gaussian noise masking should be used only for initial feasibility analysis, as its additive nature and inherent predictability make it highly sensitive to statistical analysis. The 16-bit XOR is suitable for environments with highly constrained computational resources rather than for environments lacking key management.



\section*{Acknowledgment}

The preferred spelling of the word ``acknowledgment'' in America is without 
an ``e'' after the ``g''. Avoid the stilted expression ``one of us (R. B. 
G.) thanks $\ldots$''. Instead, try ``R. B. G. thanks$\ldots$''. Put sponsor 
acknowledgments in the unnumbered footnote on the first page.

\section*{References}

Please number citations consecutively within brackets \cite{b1}. The 
sentence punctuation follows the bracket \cite{b2}. Refer simply to the reference 
number, as in \cite{b3}---do not use ``Ref. \cite{b3}'' or ``reference \cite{b3}'' except at 
the beginning of a sentence: ``Reference \cite{b3} was the first $\ldots$''

Number footnotes separately in superscripts. Place the actual footnote at 
the bottom of the column in which it was cited. Do not put footnotes in the 
abstract or reference list. Use letters for table footnotes.

Unless there are six authors or more give all authors' names; do not use 
``et al.''. Papers that have not been published, even if they have been 
submitted for publication, should be cited as ``unpublished'' \cite{b4}. Papers 
that have been accepted for publication should be cited as ``in press'' \cite{b5}. 
Capitalize only the first word in a paper title, except for proper nouns and 
element symbols.

For papers published in translation journals, please give the English 
citation first, followed by the original foreign-language citation \cite{b6}.

\begin{thebibliography}{00}
\bibitem{b1} G. Eason, B. Noble, and I. N. Sneddon, ``On certain integrals of Lipschitz-Hankel type involving products of Bessel functions,'' Phil. Trans. Roy. Soc. London, vol. A247, pp. 529--551, April 1955.
\bibitem{b2} J. Clerk Maxwell, A Treatise on Electricity and Magnetism, 3rd ed., vol. 2. Oxford: Clarendon, 1892, pp.68--73.
\bibitem{b3} I. S. Jacobs and C. P. Bean, ``Fine particles, thin films and exchange anisotropy,'' in Magnetism, vol. III, G. T. Rado and H. Suhl, Eds. New York: Academic, 1963, pp. 271--350.
\bibitem{b4} K. Elissa, ``Title of paper if known,'' unpublished.
\bibitem{b5} R. Nicole, ``Title of paper with only first word capitalized,'' J. Name Stand. Abbrev., in press.
\bibitem{b6} Y. Yorozu, M. Hirano, K. Oka, and Y. Tagawa, ``Electron spectroscopy studies on magneto-optical media and plastic substrate interface,'' IEEE Transl. J. Magn. Japan, vol. 2, pp. 740--741, August 1987 [Digests 9th Annual Conf. Magnetics Japan, p. 301, 1982].
\bibitem{b7} M. Young, The Technical Writer's Handbook. Mill Valley, CA: University Science, 1989.
\bibitem{b8} D. P. Kingma and M. Welling, ``Auto-encoding variational Bayes,'' 2013, arXiv:1312.6114. [Online]. Available: https://arxiv.org/abs/1312.6114
\bibitem{b9} S. Liu, ``Wi-Fi Energy Detection Testbed (12MTC),'' 2023, gitHub repository. [Online]. Available: https://github.com/liustone99/Wi-Fi-Energy-Detection-Testbed-12MTC
\bibitem{b10} ``Treatment episode data set: discharges (TEDS-D): concatenated, 2006 to 2009.'' U.S. Department of Health and Human Services, Substance Abuse and Mental Health Services Administration, Office of Applied Studies, August, 2013, DOI:10.3886/ICPSR30122.v2
\bibitem{b11} K. Eves and J. Valasek, ``Adaptive control for singularly perturbed systems examples,'' Code Ocean, Aug. 2023. [Online]. Available: https://codeocean.com/capsule/4989235/tree
\end{thebibliography}

\vspace{12pt}
\color{red}
IEEE conference templates contain guidance text for composing and formatting conference papers. Please ensure that all template text is removed from your conference paper prior to submission to the conference. Failure to remove the template text from your paper may result in your paper not being published.

\end{document}
